\section*{Hoofdstuk \ref{ch:b3:green-industrial} --- Groene chemie en industriële processen}

\begin{solution}{exo:b3:green-industrial:1}
Diels-Alder: 100\,\% (een additie). Verestering: $88.0/(60.0 + 46.0) = 83$\,\%. Wittig:
$104.0/(106.0 + 276.0) = 27$\,\%, waarbij trifenylfosfineoxide het grootste deel van de massa meeneemt.
\end{solution}

\begin{solution}{exo:b3:green-industrial:2}
$\mathrm{PMI} = 120/8.0 = 15$; $E = \mathrm{PMI} - 1 = 14$.
\end{solution}

\begin{solution}{exo:b3:green-industrial:3}
Eén wasbeurt van een standaardlading bij een gegeven temperatuur en
wasresultaat. Wasmiddelen worden verschillend gedoseerd: een
geconcentreerd poeder vraagt minder per wasbeurt, zodat gelijke massa's
niet dezelfde dienst leveren.
\end{solution}

\begin{solution}{exo:b3:green-industrial:4}
2-Methyltetrahydrofuraan of ethylacetaat voor dichloormethaan; heptaan
(met ethylacetaat) voor hexaan; voor DMF ethylacetaat,
2-methyltetrahydrofuraan of dimethylcarbonaat, of een koppeling in water
met oppervlakteactieve stoffen, afhankelijk van de reagentia.
\end{solution}

\begin{solution}{exo:b3:green-industrial:5}
1.000 mol zuur; 0.800 mol ester (\qty{88.0}{g/mol}): rendement 80\,\%. AE 83.0\,\%. SF $= 152/106 = 1.43$.
$\mathrm{RME} = 70.4/152 = 0.463$; $\mathrm{AE} \times \text{rendement}/\mathrm{SF} = 0.830 \times 0.80/1.43 = 0.463$.
\end{solution}

\begin{solution}{exo:b3:green-industrial:6}
Chloorhydrine: $44.0/(28.0 + 71.0 + 74.1) = 25$\,\%; directe oxidatie: 100\,\%. Eén mol $\ce{CaCl2}$ per mol
ethyleenoxide: $\qty{1.0e6}{g}/\qty{44.0}{g/mol} \times \qty{111.1}{g/mol} = \qty{2.5}{t}$ per ton.
\end{solution}

\begin{solution}{exo:b3:green-industrial:7}
$3/8 \times \qty{5.88e4}{mol} \times \qty{44.0}{g/mol} = \qty{0.97}{t}$ $\ce{CO2}$ per ton ammoniak. Voor waterstof, $\ce{CH4 + 2H2O -> CO2 + 4H2}$:
$1/4$ mol per mol, $\qty{5.0e5}{mol}/4 \times \qty{44.0}{g/mol} = \qty{5.5}{t}$ $\ce{CO2}$ per ton waterstof.
\end{solution}

\begin{solution}{exo:b3:green-industrial:8}
% ledger: dfg:H2O_l, dfh:H2O-l
Eén kilogram is \qty{500}{mol}. Uit $\Delta_rG^\circ = \qty{237.1}{kJ/mol}$: $\qty{1.19e5}{kJ} = \qty{32.9}{kWh}$. Uit $\Delta_rH^\circ = \qty{285.8}{kJ/mol}$: \qty{39.7}{kWh}
(het verschil is warmte die de omgeving kan leveren).
\end{solution}

\begin{solution}{exo:b3:green-industrial:9}
$\qty{1.0e6}{g}/\qty{146.0}{g/mol} = \qty{6.85e3}{mol}$ $\ce{N2O}$, $\qty{0.30}{t}$; $\times 273 = \qty{82}{t}$ $\ce{CO2}$-equivalent per ton adipinezuur.
\end{solution}

\begin{solution}{exo:b3:green-industrial:10}
De RME van een stap telt de overmaat reagens als verbruikt; als de
overmaat teruggewonnen en teruggevoerd wordt, is alleen het deel dat
werkelijk verloren gaat verbruikt, en is de totale massa-efficiëntie
hoger dan het product van de waarden van de stappen. In de \omterm{def:b3:green-industrial:e-factor}{E-factor} is
de gerecycleerde stroom geen afval, op voorwaarde dat de recycling
binnen de grens valt; haar energie en verliezen moeten dan meegeteld
worden.
\end{solution}

\begin{solution}{exo:b3:green-industrial:11}
Vooraf: $20 \times 0.80 = \qty{16}{kg}$ oplosmiddel, 10\,\% verloren: \qty{1.6}{kg} afval. Achteraf: \qty{4.0}{kg}, \qty{0.40}{kg} verloren. De
volledige \omterm{def:b3:green-industrial:e-factor}{E-factor} daalt met 1.2 kg per kilogram product (en de
destillatie-energie met drie vierde).
\end{solution}

\begin{solution}{exo:b3:green-industrial:12}
Per \omterm{def:b3:green-industrial:lca}{functionele eenheid} en met dezelfde grens zou ze elke
effectcategorie afzonderlijk rapporteren (klimaat, landgebruik,
watergebruik\dots) in plaats van één score. Een beslissing vereist de
gewichten die aan deze categorieën gegeven worden, de lokale schaarste
aan water en land, de onzekerheid van de gegevens, en of de biogebaseerde
grondstof met voedsel concurreert.
\end{solution}

\begin{solution}{pb:b3:green-industrial:1}
\textbf{1.} $\ce{C10H14 + C4H6O3 -> C12H16O + C2H4O2}$ (HF); $\ce{C12H16O + H2 -> C12H18O}$ (raneynikkel);
$\ce{C12H18O + CO -> C13H18O2}$ (palladium).

\textbf{2.} $\ce{C10H14 + C4H6O3 + H2 + CO -> C13H18O2 + C2H4O2}$.

\textbf{3.} $134.0 + 102.0 + 2.0 + 28.0 = \qty{266.0}{g/mol}$ beginstoffen; ibuprofen \qty{206.0}{g/mol}: AE $= 206.0/266.0 = 77$\,\%.

\textbf{4.} Isobutylbenzeen, azijnzuuranhydride, ethylchlooracetaat,
natriumethoxide, waterig zuur ($\ce{H3O+}$), hydroxylamine en water.

\textbf{5.} $134.0 + 102.0 + 122.5 + 68.0 + 19.0 + 33.0 + 2 \times 18.0 = \qty{514.5}{g/mol}$: AE $= 206.0/514.5 = 40$\,\%.

\textbf{6.} Azijnzuur, natriumchloride, ethanol, koolstofdioxide, water,
ammoniak (als ammoniumzout), en aluminiumzouten uit het
stoichiometrische aluminiumchloride.

\textbf{7.} $0.71 + 0.55 + 0.010 + 0.14 + 0.05 = \qty{1.46}{t}$.

\textbf{8.} PMI 1.46; $E = 0.46$.

\textbf{9.} $0.46 - 0.31 = 0.15$.

\textbf{10.} PMI $= 2.0 + 1.5 = 3.5$; $E = 2.5$.

\textbf{11.} Rendementen onder 100\,\%, overmaat reagentia, verliezen aan
oplosmiddel en katalysator, en opwerkingsmaterialen voegen allemaal afval
toe dat de \omterm{def:b3:green-industrial:atom-economy}{atoomeconomie} negeert.

\textbf{12.} Afval voorkomen, de \omterm{def:b3:green-industrial:atom-economy}{atoomeconomie} maximaliseren, katalysatoren
gebruiken in plaats van stoichiometrische reagentia, onnodige derivaten
en stappen vermijden, energie besparen.

\textbf{13.} Het stoichiometrische aluminiumchloride, dat bij de opwerking
tot aluminiumafval gehydrolyseerd wordt; waterstoffluoride is tegelijk
katalysator en oplosmiddel en wordt teruggewonnen en hergebruikt.

\textbf{14.} Raneynikkel (een heterogene hydrogeneringskatalysator).

\textbf{15.} Een palladiumfosfinecomplex: de carbonyleringschemie van
\cref{ch:b3:homogeneous-catalysis}.

\textbf{16.} Azijnzuuranhydride draagt één acetylgroep over; de andere
helft van het molecuul vertrekt als azijnzuur.

\textbf{17.} Waterstoffluoride is zeer giftig en bijtend, koolstofmonoxide
giftig en ontvlambaar: gesloten apparatuur uit bestendige materialen,
detectoren en opgeleide operatoren.

\textbf{18.} Elke stap voegt verwarmen, koelen, scheidingen en terugwinning
van oplosmiddel toe; drie stappen vragen van dat alles minder.

\textbf{19.} $2.5 \times 3000 = \qty{7500}{t}$ per jaar.

\textbf{20.} $0.15 \times 3000 = \qty{450}{t}$ per jaar.

\textbf{21.} Ongeveer \qty{7000}{t} afval per jaar vermeden.

\textbf{22.} Nee: de aard van het afval telt (een ton zout water is geen
ton van een persistente giftige verbinding), en ook de energie, de
emissies en de gevaren van de reagentia tellen.

\textbf{23.} De effecten van het maken van de reagentia en de energie
stroomopwaarts, van de emissies naar lucht en water, en van het gebruik
en de verwijdering van het product, per \omterm{def:b3:green-industrial:lca}{functionele eenheid}, in meerdere
effectcategorieën.

\textbf{24.} De katalytische route van drie stappen heeft een
\omterm{def:b3:green-industrial:atom-economy}{atoomeconomie} van ongeveer 77\,\% (40\,\% voor de route van zes stappen).
\end{solution}
